\section{Discussion}

We reflect on our results and relate them to other work that also analyzes annotation artifacts in NLP datasets, drawing three main conclusions.

\paragraph{Many datasets contain annotation artifacts.}
\citet{Lai2014-cd} demonstrated that lexical features such as the presence of negation, word overlap, and hypernym relations are highly predictive of entailment classes in the SICK dataset \cite{Marelli:2014}. \citet{Chen:2016} revealed problems with the CNN/DailyMail dataset \cite{Hermann:2015} which resulted from applying automatic tools for annotation. 
\citet{Levy:2016} showed that a relation inference benchmark \cite{Zeichner:2012} is severely biased towards distributional methods, since it was created using DIRT \cite{Lin:2001}.
\citet{Schwartz:2017} and \citet{Cai:2017} showed that certain biases are prevalent in the ROC stories cloze task \cite{Mostafazadeh:2016},
which allow models trained on the endings alone, and not the story prefix, to yield state-of-the-art results. \citet{Rudinger2017SocialBI} revealed that elicited hypotheses in SNLI contain evidence of various gender, racial, religious, and aged-based stereotypes.  In parallel to this work, \citet{hypothesis-only-nli} uncovered similar annotation biases across multiple NLI datasets. Indeed, annotation artifacts are not unique to the NLI datasets, and the danger of such biases should be carefully considered when annotating new datasets.

\paragraph{Supervised models leverage annotation artifacts.}
 \citet{Levy:2015} demonstrated that supervised lexical inference models rely heavily on artifacts in the datasets, particularly the tendency of some words to serve as prototypical hypernyms. \citet{Agrawal:16,Jabri:2016,Goyal:17} all showed that state-of-the-art visual question answering~\cite{Antol:15} systems leverage annotation biases in the dataset. \citet{Cirik:18} find that complex models for referring expression recognition achieve high performance without any text input. In parallel to this work, \citet{Dasgupta2018-ry} found that the InferSent model \cite{conneau2017supervised} relies on word-level heuristics to achieve state-of-the-art performance on SNLI. These findings coincide with ours, and strongly suggest that supervised models will exploit shortcuts in the data for gaming the benchmark, if such exist.

% \paragraph{Annotation artifacts mask the model's actual performance level.}
\paragraph{Annotation artifacts inflate model performance.}
This is a corollary of the above, since large portions of the test set can be solved by relying on annotation artifacts alone. 
A similar finding by \citet{Jia:2017} showed that the performance of top question-answering models trained on SQuAD \cite{Rajpurkar:2016} drops drastically by introducing simple adversarial sentences in the evidence. 
We release the \hard SNLI and MultiNLI test sets,\footnote{SNLI: \url{goo.gl/5rQKb5}, MultiNLI matched: \url{goo.gl/abdSbi}, MultiNLI mismatched: \url{goo.gl/Cu9Gp6}\\}
and encourage the community to use them for evaluating NLI models (in addition to the original benchmarks).
We also encourage the development of additional challenging benchmarks that expose the true performance levels of state-of-the-art NLI models.

